\section{Scope and Limitations}

The results in this paper are architecture-level claims, not a substitute for a full benchmark campaign
on a production agent stack. Section~6 compares the topology results to external architecture-level
aggregates from Kim et al.~\cite{kim2025scaling}, but that comparison is qualitative. It does not fit the
theorem parameters directly, and it is not an end-to-end evaluation of the Operon runtime on the same
benchmark suite.

The executable substrate described in Section~7 is likewise narrower than a full semantic verification
framework. The finite-trace comparison machinery operates on supplied deterministic input scripts, so it
does not establish general behavioral equivalence for nondeterministic, tool-rich, or continuously
interactive agent systems. The implementation shows that the topology model can be instantiated and
inspected, not that all possible rewritings or coordination patterns are formally certified.

Finally, the topology classes studied here are intentionally coarse. Real deployments often mix reviewer
gates, hub coordination, local feedback, and dynamic tool routing inside the same workflow. Model quality,
prompt design, and tool reliability can also dominate topology effects in specific domains. The present
paper should therefore be read as a compact architectural calculus for agent coordination, with direct
benchmarking and adaptive topology selection left to future work.  A companion paper~\cite{banu2026benchmarks}
provides initial end-to-end evaluation of Operon's structural guarantees with real LLM agents, partially
addressing the system-level evaluation gap noted above.
